\chapter*{Abstract}
\addcontentsline{toc}{chapter}{Abstract}
Decisions in retail investing are influenced not only by movements in the market but also by unrealized gains and losses, the liquidity of the household, the holding period and the reaction of the investor to uncertainty. Tools that predict prices can hardly tell whether a particular intended decision is the result of a behavioral pattern, or what the alternative course of action would mean for the investor's financial position in the future. In this project we propose Folium, an integrated decision support system and web application that evaluates one intended action (sell, hold or buy more) on one holding against the whole financial picture of the investor.

For every analysis, Folium reads the reason the investor gives in their own words for any stated need for cash, and then classifies the most likely behavioral pattern behind the decision (loss aversion, herd behavior, the disposition effect or a rational motive) using a calibrated XGBoost classifier with 22 features and exact TreeSHAP attributions. It compares the wealth of holding and of selling over the investor's own horizon on the same 1,000 block-bootstrapped market paths built from real NIFTY 50 and gold history, after Indian capital gains tax and charges. Cause and effect is estimated with DoWhy backdoor adjustment and g-computation. A fact-driven template engine explains the numbers and assumptions in English and Hindi, the verdict (HOLD, SELL or BUY MORE) comes from written evidence rules rather than a language model, and every analysis is sealed in a per-investor SHA-256 hash chain that can be verified later. The system is built with FastAPI, PostgreSQL with row-level security, Redis and React.

On a held-out test set of 4,000 synthetic decisions placed on real market days, the classifier achieves an accuracy of 87.9\%, a macro-F1 score of 0.840 and an expected calibration error of 0.014. It is significantly better than a rules-only baseline and logistic regression, and statistically tied with a random forest. Among six market path generators, only the 12-month block bootstrap reproduces both the spread of real one-year returns and volatility clustering. The causal estimates recover the known generator effects with a mean absolute error of about 0.5 percentage points. These results are at the development stage and are obtained using synthetic behavioral labels instead of real investors' decisions.

\textbf{Keywords:} behavioral finance, decision support, XGBoost, probability calibration, SHAP, causal inference, block bootstrap, Monte Carlo simulation, explainable AI.

